\section{Panel Overview}

The empirical analysis uses the replication package supplied with this
thesis rather than re-downloading every raw input from its original upstream
provider.
The base file, \path{crypto_panel_features.npz}, is a processed weekly
cryptocurrency panel with $T = 324$ calendar weeks
(2020-01-05 to 2026-03-15), $N = 86$ tickers, one next-week return target,
and 49 pre-existing predictor columns.
The augmented file, \path{crypto_panel_features_with_etf.npz}, appends
four granular ETF-flow variables, so the final research panel has 53
predictor columns.
The six information sets used in the empirical comparisons draw on a
42-feature subset of those columns.

The base panel is treated as a processed input because the current repository
does not contain enough raw-vendor metadata to reconstruct every original
source for the 49 pre-existing variables.
For that reason, the data discussion below separates directly observable
replication files from upstream provenance that is not recoverable from the
files provided.
This is a deliberate reporting constraint: feature names and transformations
are reported when they can be inspected in the files and code, but external
vendors are not attributed unless the repository or an official source
supports the attribution.

Missing values are encoded as $-99.99$ and masked in both loss computation
and portfolio evaluation.
The fixed 86-token universe is therefore best interpreted as the asset
universe encoded in the replication package, not as a fully auditable
historical listing universe.

\section{Observed Data Files and Provenance}

\Cref{tab:data_provenance} lists the data files that can be directly verified
in the replication package.
Only the two Farside CSV files are raw external downloads visible in the
repository.
The NPZ files are processed research inputs or saved benchmark arrays.

\begin{table}[t]
\centering
\caption{Directly observed data files in the replication package.}
\label{tab:data_provenance}
\scriptsize
\begin{tabularx}{\linewidth}{>{\raggedright\arraybackslash}p{0.20\linewidth}X >{\raggedright\arraybackslash}p{0.16\linewidth}X >{\raggedright\arraybackslash}p{0.12\linewidth}}
\toprule
\thead{File} & \thead{Variables / role} & \thead{Period} &
\thead{Construction step} & \thead{Raw source observable?} \\
\midrule
\path{crypto_panel_features.npz}
  & Processed base panel: dates, tickers, next-week returns, and 49
    pre-existing predictors.
  & Weekly, 2020-01-05--2026-03-15
  & Loaded as the pre-existing base panel; no raw-source reconstruction is
    possible from the current files.
  & No \\
\path{crypto_panel_features_with_etf.npz}
  & Final v6 panel after appending four Farside ETF variables.
  & Weekly, 2020-01-05--2026-03-15
  & Constructed by augmenting the base panel with weekly GBTC, BTC ex-GBTC,
    ETHE, and ETH ex-ETHE flow variables.
  & Partly \\
\path{market_cap.npz}
  & Saved market-cap matrix used for market benchmark construction.
  & Weekly, 2020-01-05--2026-03-15
  & Loaded as an aligned matrix for weighting and benchmark calculations.
  & No \\
\path{farside_btc.csv}
  & Daily U.S. spot Bitcoin ETF flow table, including GBTC and total flows.
  & Daily rows, 2024-01-11--2026-03-06
  & Parsed from CSV, converted to signed values, then aggregated to weekly
    GBTC and BTC ex-GBTC flows.
  & Yes \\
\path{farside_eth.csv}
  & Daily U.S. spot Ethereum ETF flow table, including ETHE and total flows.
  & Daily rows, 2024-07-23--2026-03-06
  & Parsed from CSV, converted to signed values, then aggregated to weekly
    ETHE and ETH ex-ETHE flows.
  & Yes \\
\bottomrule
\end{tabularx}
\end{table}

The Farside files correspond to the Bitcoin and Ethereum ETF flow tables
published by Farside Investors \citep{bitcoin_etf_farside,ethereum_etf_farside}.
In the CSV format, parenthesized values denote outflows and are parsed as
negative numbers, dashes are treated as missing or zero-flow cells depending
on the aggregation step, and daily rows are aggregated to the weekly panel
calendar.
For each week, flows in the interval from the previous panel date to the
current panel date are summed.
The four appended variables are GBTC flow, total Bitcoin ETF flow excluding
GBTC, ETHE flow, and total Ethereum ETF flow excluding ETHE.
Each weekly series is then transformed with the same 52-week rolling
z-score rule used for other time-series predictors.

\section{Feature Categories}

\Cref{tab:features} summarizes the 42 features actually used in the model
comparisons.
The table reports variable groups and examples; it should not be read as a
claim about raw upstream vendors for the pre-existing base-panel variables.
The grouping reflects the main conceptual distinction in the paper: some
signals are slow-moving and structural, whereas others are event-like and
time-dependent.

\begin{table}[t]
\centering
\caption{Feature categories used in the empirical model comparisons.
         $M_g$ = number of features in group $g$.}
\label{tab:features}
\small
\begin{tabular}{llc}
\toprule
\thead{Group} & \thead{Examples} & \thead{$M_g$} \\
\midrule
Price \& Momentum    & 1w/4w/12w/26w/52w returns               & 5  \\
Technical            & RSI, Bollinger \%, ATR, OBV, volume ratio & 6  \\
On-chain             & Active addresses, NVT, MVRV, net flow   & 5  \\
Macro \& Sentiment   & Fear-greed index, equity, FX, volatility, metal returns & 11 \\
ETF \& Prediction Market & BTC/ETH ETF flow fields, BTC probability field & 6  \\
Trump Social Media   & Post count, caps ratio, tariff score, crypto score, sentiment & 5  \\
Farside ETF Granular & GBTC, BTC ex-GBTC, ETHE, ETH ex-ETHE    & 4  \\
\midrule
Total                &                                          & 42 \\
\bottomrule
\end{tabular}
\end{table}

\section{Data Splits}
\label{sec:data:splits}

We partition the panel chronologically into training (70\%, weeks 0--226),
validation (15\%, weeks 227--274), and test (15\%, weeks 275--323),
corresponding to 49 held-out panel weeks.
Because the saved prediction masks exclude one test week from valid
long-short portfolio construction, the reported Sharpe ratios are computed
from 48 weekly portfolio-return observations.
Hyperparameters are selected on the validation set, and the test set is used
only for final reporting.
This preserves the temporal ordering of the problem and avoids contaminating
evaluation with repeated test-period feedback.

\section{Look-ahead Bias}
\label{sec:data:lookahead}

The replication code applies time-series standardization using a rolling
52-week z-score.
At week $t$, the mean and standard deviation are computed over information
available through week $t$.
No future panel dates enter the normalization window.

For the Farside ETF variables, daily rows are first assigned to the
corresponding panel week and then standardized.
The inclusion of week $t$ in its own rolling mean gives the current
observation weight $1/52 \approx 2\%$ once the full window is available.
This avoids future leakage, although it still means that weekly ETF-flow
features are contemporaneous weekly aggregates rather than signals observed
strictly before the start of the week.
The same timing interpretation applies to all processed weekly predictors
whose raw timestamp conventions are not recoverable from the base panel.

\section{Provenance and Survivorship Limitations}
\label{sec:data:survival}

Two limitations follow from the available replication files.
First, the upstream raw sources for the 49 pre-existing base-panel variables
cannot be fully audited from the current repository.
This affects source attribution, not the internal consistency of the model
comparisons: every method uses the same processed panel, feature
configurations, masks, and chronological splits.

Second, the fixed 86-token universe may contain survivorship bias.
The repository provides the final ticker list but not a complete historical
listing-and-delisting rule.
Tokens that were historically relevant but absent from the fixed universe
would therefore be excluded from the analysis.
The paper should consequently be read primarily as evidence on the relative
suitability of objective functions and model architectures for a fixed
cross-sectional cryptocurrency ranking problem, rather than as an
unconditional claim about live investability of the entire cryptocurrency
market.
